\chapter{Strata Forth}\label{ch:forth}

Strata Forth (\texttt{lang/forth/}) is an interactive programming system for the bare machine: a small
kernel in assembly language (\texttt{forth.s}), and the rest of the language written in Forth
(\texttt{core.fs}) and compiled at power-on. It is the manual's example of a language implementation
that takes the architecture's rules as its design: words compile to native code that reaches the
instruction Local Store the only way code can, the storage tiers and capabilities are words, and the
other cores are workers that are lent the image.

\begin{shell}
make forth                 the REPL on the emulator, on every core
make forth-rtl             the same on the RTL simulator
make forth-demo-tiers      one demo: tiers, multicore or gfx; make forth-demos runs all three
make test-forth            the regression test on both machines
\end{shell}

Input may be typed or fed from a file with \texttt{-i}; each line is echoed, then its output follows,
then \texttt{ok}. The end of input, or \texttt{bye}, halts every core. The REPL runs in supervisor mode,
which nothing can preempt, so there is no Ctrl-C.

\begin{shell}
Strata Forth 0.1 (native code, 8 cores, bare metal) 32528 code bytes free, 18200 data bytes free
: sq ( n -- n*n ) dup * ;  12 sq . 144  ok
true jit-trace  : cube dup dup * * ; [jit: tput+tgetx 64 bytes -> I-LS 0x8160]  ok
1234 0 far!  0 far@ . 1234  ok
.fcache 1 hits, 1 misses, 0 write-backs
7 1 ' sq 1 spawn  1 await . 49  ok
\end{shell}

\section{Native code, installed by transfers}\label{sec:forth-jit}

A colon definition is a machine-code subroutine: a prologue that saves \reg{ra}, one \insn{jal} per
word called or the word's body copied inline, and an epilogue.

\begin{center}\small
\begin{tabular}{@{}l l@{}}
\toprule
\textbf{Register} & \textbf{Use}\\\midrule
\reg{s1} & the top of the data stack, cached\\
\reg{s0} & the data stack pointer: it points at the second item and grows down\\
\reg{sp} & the return stack, which is the machine stack\\
\reg{s2}, \reg{s3} & the index and limit of the innermost \texttt{do} loop\\
\reg{s4} & the locals frame\\
\reg{t0}--\reg{t6}, \reg{a0}--\reg{a7}, \reg{s5}--\reg{s11} & scratch\\
\bottomrule
\end{tabular}
\end{center}

Stores cannot reach the I-LS, and the I-LS cannot be read. Strata Forth therefore keeps an
\term{identity mirror} of its I-LS in far memory: far offset = I-LS address. The compiler works in three
steps.
\begin{enumerate}
\item \texttt{:} opens a definition at the next free I-LS address. Instructions are assembled into a
  4~KiB staging buffer in the D-LS, where forward branches are patched.
\item \texttt{;} compiles the epilogue and installs the code with two transfers:
  \insn{tput} from the staging buffer to the mirror, then \insn{tgetx} from the mirror into the I-LS at
  the same address.
\item The word can run at once: a fetch from a range that a \insn{tgetx} is still filling waits
  (Section~\ref{sec:asif}).
\end{enumerate}

The kernel installs itself the same way: the image's \texttt{.text} is only a boot stub, which
\insn{tgetx}es the kernel from the far image. Everything compiled is therefore readable in far memory.
\texttt{see} disassembles a word by reading it from the mirror; the compiler inlines a primitive by
reading its body from the mirror with \insn{tget}; and other cores install the code from there.

\begin{shell}
see sq
: sq  ( xt $80f0 , 40 bytes)
  80f0:   442fff8  addi sp, sp, -8
  80f4:  60220000  sd ra, 0(sp)
  80f8:   508fff8  addi s0, s0, -8
  80fc:  61280000  sd s1, 0(s0)
  8100:  40a80000  ld t0, 0(s0)
  8104:   5080008  addi s0, s0, 8
  8108:   125480a  mul s1, t0, s1
  810c:  40220000  ld ra, 0(sp)
  8110:   4420008  addi sp, sp, 8
  8114:  a4010000  ret
\end{shell}

In \texttt{sq}, the bodies of \texttt{dup} and \texttt{*} were copied inline between the prologue and the
epilogue.

The compiler inlines primitives of up to four instructions, folds a 16-bit literal into the
instruction that uses it (a variable followed by \texttt{@} is one \insn{ld} from an absolute D-LS
address), turns a call before \texttt{;} into a jump, and keeps \texttt{do} loop indices in registers.
Control structures, loops, locals and quotations are themselves written in Forth with the
code-generation words. Code space and data space are separate, as the Local Stores are: \texttt{here},
\texttt{,} and \texttt{allot} address the dictionary in the D-LS, and \texttt{chere} is the code pointer.

\subsection{Instructions built at run time}
\IS instructions name capability registers and transfer tags in immediate fields, so ``transfer
through capability register $n$'' is not one instruction but sixteen. \texttt{xexec} assembles the
instruction wanted, installs it followed by \insn{ret} into one of eight I-LS trampolines (a
\insn{tput} and a \insn{tgetx} of 8 bytes; the trampoline is chosen by a hash of the instruction, so a
loop pays the install once) and calls it. Every capability and transfer word is a one-line definition
over \texttt{xexec}.

\section{Words for the architecture}

\begin{table}[htbp]
\centering\small
\caption{Strata Forth words for tiers, capabilities and cores. \textit{c\#} is a capability register
number; a transfer through register $n$ uses tag $n$.}\label{tab:forth}
\begin{tabularx}{\linewidth}{@{}>{\ttfamily\raggedright\arraybackslash}p{15.5em} X@{}}
\toprule
\textrm{\bfseries Words} & \textbf{Function}\\\midrule
tget tput ( dls off len c\# -- ) & a transfer between the D-LS and the segment of a capability\\
tag ( c\# -- mask ) \ twait tpoll tany & tags: wait for, poll, wait for any\\
csplit ( cs off cd -- ) \ cjoin ( cd cs -- ) & split and join segments\\
cmov cshare catomic ( cs cd -- ) & move, share, make atomic\\
crestrict ( c\# perms -- ) \ cclear ( c\# -- ) & restrict, discard\\
cbase clen cperms cvalid ( c\# -- x ) & inspect; \texttt{.cap} and \texttt{.caps} print capabilities\\
give take ( c\# core -- ) & \insn{cgive} to a core; \insn{ctake} from a core\\
tamo ( addr operand c\# fn -- old ) & a far atomic; \texttt{atomic@}, \texttt{atomic!}, \texttt{atomic+!}\\
far@ far! farc@ farc! & cells and bytes of the far heap, through the software cache\\
far-read far-write ( dls far u -- ) & the far heap in bulk\\
fflush finval .fcache fcache-reset & write back, forget, show and reset the cache's counters\\
cores core\# workers & the number of cores, this core, the number of workers\\
spawn ( x1..xn n xt core -- ) & start a word on another core with $n$ arguments\\
await ( core -- y1..ym ) & wait for it and fetch its results\\
run-on ( xt core -- ... ) \ par ( xt -- ) & spawn and await; run a word on every worker\\
ticks ( -- us ) \ ms ( n -- ) & the timer's NOW; sleep on the timer's WAIT\\
cycles instret timeit ( xt -- ) csr@ & cost: counters, the time and cycles of a word, a CSR\\
see disasm dis code@ & the disassembler; code read back from the mirror\\
.jit jit-trace ( f -- ) chere & the code space, a report of every install\\
xexec ( a0 a1 a2 insn -- a0' ) & run one instruction built at run time\\
insn, mark if, ahead, resolve addi, ld, \ldots & code generation and the assembler\\
\bottomrule
\end{tabularx}
\end{table}

\subsection{Far memory through a software cache}
\texttt{far@} and \texttt{far!} address the heap capability through a 2-way, 32-set write-back cache of
64-byte lines in the D-LS, built on \insn{ltag} and \insn{linst} (Section~\ref{sec:swcache}). A miss
writes a dirty victim back with \insn{tput} and fetches the line with \insn{tget}; the load from the
line simply waits for the transfer. \texttt{.fcache} shows hits, misses and write-backs, so the cost of
a far reference is on display. Each core has its own cache, over whatever its heap register holds.

\subsection{Cores as workers}
At power-on core 0 gives each worker a private 64~KiB page and a copy of the \term{job board}, a 4~KiB
atomic segment of the cluster store. A worker waits in \insn{ctake}.
\begin{itemize}
\item When it receives a capability with \perm{X}, that is the \emph{image}: core 0's own memory
  capability, lent. The worker \insn{tgetx}es the compiled code into its I-LS, \insn{tget}s a snapshot
  of the dictionary into its D-LS, and gives the capability back. Code compiled on core 0 thus runs on
  every core, although no core can write another's I-LS.
\item Otherwise it is the \emph{doorbell}: the worker reads its job slot on the board with far atomics,
  runs the word under \texttt{catch}, writes back its results and rings core 0.
\end{itemize}
\texttt{spawn} publishes a new image only when something was compiled since the last one, and lends it
only to workers that hold an older generation. Memory moves between cores as ownership:

\begin{shell}
: job ( k -- sum )  drop  8 0 take              \ our slice of the heap becomes our c8
  0  n 0 do  i cells far@ +  loop               \ sum it through our own far cache
  8 0 give ;                                    \ hand it back
: deal ( -- )  cores 1 ?do  i 1 ['] job i spawn  loop
  cores 1 ?do  8 slice 11 csplit  8 i give  11 8 cmov  loop ;
\end{shell}

The three demos in \texttt{lang/forth/demos/} are sessions of this kind: \texttt{tiers.fs} walks
through capabilities, transfers, the far cache, a bounds fault caught with \texttt{catch}, ownership
moving to another core and back, and code installation; \texttt{multicore.fs} is the parallel sum
above; \texttt{gfx.fs} cuts the framebuffer into one band per worker, renders the Mandelbrot set in
parallel, joins the bands and writes PRESENT.

\section{Memory map and dialect}

\begin{center}\small
\begin{tabularx}{\linewidth}{@{}l l X@{}}
\toprule
\textbf{Where} & \textbf{Range} & \textbf{Contents}\\\midrule
I-LS & \hex{0000}--\hex{013F} & the boot stub, the worker loop, eight \texttt{xexec} trampolines\\
I-LS & \hex{0140}-- & the kernel (about 11~KiB), then compiled code; about 32~KiB free after start-up\\
D-LS & \hex{0000}-- & kernel variables and buffers, the far cache, then the dictionary (about 18~KiB free)\\
D-LS & \hex{9800}--\hex{9EFF} & the data stack\\
D-LS & \hex{E000}--\hex{FFFF} & the return stack\\
T3 & $[0, 1~\text{MiB})$ & \reg{c0}: the I-LS mirror, the text of \texttt{core.fs}, the D-LS snapshot for workers\\
T3 & $[1, 2~\text{MiB})$ & the workers' private pages\\
T3 & $[2, 48~\text{MiB})$ & \reg{c8}: the heap behind \texttt{far@} and \texttt{far!}\\
T3 & $[48, 64~\text{MiB})$ & \reg{c9}: the framebuffer window\\
T2 & $[0, 4~\text{KiB})$ & \reg{c12}: the job board, atomic\\
\bottomrule
\end{tabularx}
\end{center}

The dialect is close to Forth 2012 for the core words. A cell is 64 bits and there are no double-cell
words; \texttt{/} and \texttt{mod} truncate, like the hardware's division; names are case-blind;
numbers may be written \texttt{\$ff}, \texttt{\#99}, \texttt{\%1010} and \texttt{'c'}. Beyond the
standard words it has locals (\texttt{\{ a b -- \}}), quotations (\texttt{[: \ldots\ ;]}),
\texttt{catch} and \texttt{throw} with error messages that name the word, \texttt{words}, \texttt{.s}
and \texttt{dump}. Compiling \texttt{core.fs} at power-on takes about one million instructions.

Strata Forth runs on the bare machine. A program of its kind under StrataOS, which installs code at
run time, declares its far mirror to the kernel as a code mirror, so that a context switch reloads the
I-LS from it (Section~\ref{sec:os-switch}).

\subsection{Code paging}\label{sec:forth-paging}
Under StrataOS a program has 32~KiB of I-LS. A Forth of one thread keeps only its kernel there:
16~KiB, its code mirror. Everything compiled after it, \texttt{core.fs} on, goes into \emph{code
units}: 2~KiB slot-aligned pieces of a store in far memory (104~KiB of \reg{c0}), each holding whole
definitions. StrataOS pages them (Section~\ref{sec:os-paging}): the program sets \csr{xslot} to 11,
a call of a unit that is not in the I-LS misses, and the dispatcher copies the unit into a slot of
the code cache, the I-LS after the mirror (seven slots), maps it with \insn{xmap} and retries.
Forth's own trap handler, which turns faults into \texttt{throw}s, gets the miss first and hands it
over with svc PAGE.

A unit's code runs in whatever slot it got, so the compiler makes it independent of its place. A
call of another unit, or of a word in the same one, is \texttt{addi t5, tp, \textit{offset}} and
\insn{xcall}~\texttt{ra, c0, t5}, which leaves a far link: the return pages the caller back in if it
was evicted meanwhile. A call of the kernel is \texttt{jalr ra, gp, \textit{offset}} (\reg{gp} holds
\hex{10000}), a plain link, safe because the kernel never pages, except for the few kernel words that
run code they are given (\texttt{execute}, \texttt{catch}, \texttt{evaluate}, \texttt{included},
\ldots): those are reached through their entries in a gateway unit, by \insn{xcall}, and return with
\insn{xret}. The prologue and the return go through \reg{gp} to the kernel's shared ones. An execution
token is the code's address: an I-LS address for the kernel, a far offset for a unit. A definition is
staged as before and placed at \texttt{;}: where the last one ended, or at the start of the next unit
if it does not fit there (its few absolute references to itself move with it); installing it is a
\insn{tput} to the unit, after which the overlay map is cleared, since a stale copy may be in the
cache. A definition must fit in a unit. The dispatcher puts a program's cache back as it was after
every context switch, since a slot may hold the code a plain link, or the interrupted program
counter, leads back to. With threads (\texttt{forth -n N}) every word is resident, as before: a
worker installs the code from the parcel it is lent.

The editor \texttt{lang/forth/vi.fs}, the Forth port of StrataOS's \texttt{vi}
(Section~\ref{sec:os-vi}), is some 60~KiB of code in units, with no annotations; a key costs about
what it did when the editor was built on code overlays of its own (54~K instructions). Its byte loops
are code words assembled by \texttt{core.fs}'s assembler. Data space is its limit (16~KiB holds a
chunk of text, the line being edited and every word's name), so its text, undo history, register,
keyword tables and message texts are in far memory.

\subsection{The ports of StrataOS's programs}\label{sec:forth-progs}
Twenty-two of StrataOS's programs are also written in Strata Forth, \texttt{lang/forth/progs/},
and run as \texttt{forth NAME.fs ARGS}: the same arguments, output, error messages and exit codes,
in pipes and with redirections. At the Forth prompt \texttt{include NAME.fs} defines the word
\texttt{NAME}, which runs the program on the rest of its line.

\begin{center}\small
\begin{tabular}{@{}l l@{}}
\toprule
\textbf{Ports} & \textbf{What they are}\\\midrule
\texttt{seq cat upper wc grep hexdump} & data through standard input and output, a disk block\\
\texttt{touch rm cp mv tag untag retag tags ls} & files and tags, through the file server\\
\texttt{df free uptime devs cores ps sched} & status reports from \texttt{fs}, \texttt{proc} and \texttt{dev}\\
\texttt{spin units} & a busy loop; 24 functions of 1~KiB, each in a code unit of its own\\
\bottomrule
\end{tabular}
\end{center}

They share a library, \texttt{sys.fs} (the command line, standard input by lines, standard error,
messages to the servers; about 2.9~KiB of code). The disk's directory has room for 128 files, so
the ports are packed into one, \texttt{progs.pk}: \texttt{included}, given a name that is no file,
runs the word in a kernel variable instead, and that includes the pack, whose head
(\texttt{progs.fs}) finds the port in its index and interprets that part of the pack's text. A start costs about 2.5~M instructions more
than the assembly program; past that, with the byte loops as code words, a long input takes about
1.1 to 1.8 times its instructions. \texttt{tools/test-fs-progs.sh} runs every port against its
assembly program.
